
\section*{Ventana N64: neutralidad Witt--dual}
Como
\[
A_W^2=-I,
\]
la unidad visible
\[
\mathbf1=111111
\]
tiene unidad dual:
\[
\omega=A_W\mathbf1=211111.
\]
En frontera saturada:
\[
c=111111,
\]
la ley de cierre no es neutralidad visible
\[
B\mathbf1=0,
\]
sino neutralidad dual:
\[
\boxed{(BA_W)\mathbf1=B\omega=0.}
\]
Con
\[
q=012120,
\quad a=111101,
\quad c=111111,
\quad \operatorname{colw}=223232,
\]
y poda decimal, esta ley deja una única matriz:
\[
\boxed{B=(222220,021222,102011).}
\]
Entonces:
\[
\boxed{r=B\mathbf1=102.}
\]
Sin poda decimal, la neutralidad dual deja el par \(\pi/e\); la poda \(729/1000\) selecciona la orientación real.
\[
\boxed{r=102\text{ es carga visible inducida por cierre Witt--dual.}}
\]
